\documentclass{article}
\usepackage[a4paper, margin=3cm]{geometry}
\usepackage{hyperref}

    \usepackage{enumitem}
    \newlist{list_1}{enumerate}{1}
    \setlist[list_1,1]{
        label=\roman*.
    }
    \newlist{list_2}{enumerate}{1}
    \setlist[list_2,1]{
        label=\roman*.,
        leftmargin=2cm
    }
    \newlist{list_3}{enumerate}{1}
    \setlist[list_3,1]{
        label=\roman*.,
        leftmargin=3.5cm
    }

% Package for header and footer
    \usepackage{fancyhdr}
    \pagestyle{fancy}
    \fancyhf{}
    \fancyhead[L]{Fisica}
    \fancyhead[C]{\bfseries{Analisi I}}
    \fancyhead[R]{Olmo Astorino}
    \fancyfoot[C]{\thepage}

% Package for mathematical symbols
    \usepackage{amsmath}
    \usepackage{amssymb}
    \usepackage{amsfonts}
    \usepackage{amsthm}

    \newtheoremstyle{nodot}
        {11pt}    % Space above
        {11pt}    % Space below
        {\normalfont} % Body font
        {1cm}       % Indentation
        {\bfseries} % Heading font
        { /}       % Punctuation after heading
        {0.5em}  % Space after heading
        {}       % Heading specification
    \theoremstyle{nodot}

    \newtheorem{definition}{Def}[subsection]
    \newtheorem{axiom}{Ass}[subsection]
    \newtheorem{proposition}{Pro}[subsection]
    \newtheorem{theorem}{Teo}[subsection]
    \newtheorem{lemma}{Lem}[subsection]
    \newtheorem{corollary}{Cor}[subsection]
    \newtheorem{property}{P}[subsection]
    \newtheorem{demonstration}{Dim}[subsection]

    \newtheoremstyle{obs}
        {11pt}    % Space above
        {11pt}    % Space below
        {\normalfont} % Body font
        {}       % Indentation
        {\bfseries} % Heading font
        { /}       % Punctuation after heading
        {0.5em}  % Space after heading
        {}       % Heading specification
    \theoremstyle{obs}
    \newtheorem{observation}{Oss}[subsection]

    \newtheoremstyle{examplestyle}
        {16pt}       % Spazio sopra
        {16pt}       % Spazio sotto
        {\normalfont}
        {2cm}       % Indentazione del corpo
        {\bfseries} % Stile dell'intestazione
        {:}         % Punteggiatura
        {0.5em}     % Spazio dopo l'intestazione
        {}

    \theoremstyle{examplestyle}
    \newtheorem{example}{Es}[subsection]

\begin{document}

\begin{titlepage}
    \thispagestyle{empty}
    \centering

    \vspace*{3cm}

    \huge{Fisica}
    \par

    \vspace{1.5cm}

    \Huge{\bfseries{Analisi I}}
    \par

    \vspace{1.5cm}

    \LARGE{Olmo Astorino}
    \par

    \vfill

    \large{\bfseries{UniMi - a.a. 2026/2027}}

\end{titlepage}

\tableofcontents

\newpage

\section{Insiemi numerici}

\subsection{Numeri naturali}

Definiamo i numeri naturali come segue.

\begin{definition}
    \textbf{Notazione dei numeri naturali}
    \center

    L'insieme dei numeri naturali è denotato con $\mathbb{N}$.
\end{definition}

\begin{definition}
    \textbf{Numeri naturali}
    \label{def:naturals}
    \center
    
    $ \mathbb{N}$ è definito tramite le due seguenti proprietà caratteristiche.
    \begin{list_2}
        \item $1 \in \mathbb{N}$.
        \item $\forall n \in \mathbb{N}$, $ n + 1 \in \mathbb{N}$ (proprietà del successore).
    \end{list_2}
\end{definition}

\begin{observation}
    I numeri naturali non contengono lo zero, ma è possibile estendere
    l'insieme dei numeri naturali per includere lo zero, definendo così l'insieme $\mathbb{N}_0$.
\end{observation}

\begin{definition}
    L'estensione di $\mathbb{N}$ per includere lo zero è
    \[
        \mathbb{N}_0 := \mathbb{N} \cup \{0\}.
    \]
\end{definition}
    

\vspace{0.7cm}
\subsubsection{Principio di induzione}

Da \textit{\hyperref[def:naturals]{def~\ref*{def:naturals}}} segue direttamente il principio di induzione.

Il principio di induzione è una tecnica dimostrativa tipica delle proprietà delle successioni,
che permette di dimostrare la veridicità di ogni proposizione appartenente a una famiglia
di affermazioni, la quale è in relazione biunivoca con i numeri naturali o un loro sottoinsieme.

\vspace{0.5cm}

\textbf{Principio di induzione:} Sia $\left\{P_n\right\}_{n \geq n_0}$ una famiglia di affermazioni con $n_0, n \in \mathbb{N}$. Per dimostrare la veridictà di ognuna affermazioni è sufficiente dimostrare le seguenti due proposizioni.
\begin{list_1}
    \item \textbf{Base dell'induzione:} $P_{n_0}$
    \item \textbf{Passo induttivo:} $\forall n \geq n_0$, $P_n \rightarrow P_{n+1}$
\end{list_1}

\begin{example}
    Si dimostri per induzione 
    \[
        \forall n \in \mathbb{N}, \quad \sum_{k=1}^{n} k = \frac{n(n+1)}{2}.
    \]

    \begin{list_3}
        \item \textbf{Base dell'induzione:}
        \[
            Hp: n_0=1, \qquad Th: \sum_{k=1}^{n} k = \frac{(n)(n+1)}{2}
        \]
        \vspace{0.2cm}
        \[
            Dim:
        \]
        \[
            \sum_{k=1}^{1} k = 1 = \frac{1(1+1)}{2}. \qed
        \]
        \item \textbf{Passo induttivo:}
            \[
                Hp: \sum_{k=1}^{n} k = \frac{n(n+1)}{2}, \qquad Th: \sum_{k=1}^{n+1} k = \frac{(n+1)(n+2)}{2}
            \]
            \[
                Dim:
            \]
            \[
                \sum_{k=1}^{n+1} k =
                \sum_{k=1}^{n} k + (n+1),
            \]
            Per $Hp$ si ha che
            \[
                \sum_{k=1}^{n} k + (n+1) =
                \frac{n(n+1)}{2} + (n+1),
            \]
            dunque
            \[
                \frac{n(n+1) + 2(n+1)}{2} =
                \frac{(n+1)(n+2)}{2},
            \]
            per proprietà transitiva
            \[
                \sum_{k=1}^{n+1} k =
                \frac{(n+1)(n+2)}{2}. \qed
            \]
    \end{list_3}

\end{example}

\vspace{0.7cm}
\subsubsection{Operazioni in $\mathbb{N}$}

All'interno di $\mathbb{N}_0$ si può definire l'operazione di somma "$+$", che gode delle seguenti proprietà.

\begin{property}
    \textbf{Commutativa}
    \[\forall a, b \in \mathbb{N}_0, \, \, a + b = b + a.\]
\end{property}
\begin{property}
    \textbf{Associativa}
    \[\forall a, b, c \in \mathbb{N}_0, \, \,  (a + b) + c = a + (b + c).\]
\end{property}
\begin{property}
    \textbf{Esistenza dell'elemento neutro}
    \[\forall a \in \mathbb{N}_0, \, \,  a + 0 = 0 + a = a.\]
\end{property}
\begin{property}
    \textbf{Chiusura di $\mathbb{N}_0$ rispetto alla somma} 
    \[\forall a, b \in \mathbb{N}_0, \; (a + b) \in \mathbb{N}_0.\]
\end{property}

\begin{observation}
    All'interno di $\mathbb{N}_0$ non può essere definito l'elemento inverso rispetto alla somma
\end{observation}

\subsection{Numeri interi}


Dall'estensione di $\mathbb{N}_0$ per includere gli elementi inversi rispetto alla somma,
si ottiene l'insieme dei numeri interi, denotato con $\mathbb{Z}$. Questo insieme gode delle stesse proprietà di $\mathbb{N}_0$ rispetto alla somma, con l'aggiunta della proprietà dell'esistenza dell'elemento inverso.

\begin{property}
    \textbf{Esistenza dell'elemento inverso}
    \[\forall a \in \mathbb{Z}, \, \, \exists b \in \mathbb{Z} \mid a + b = b + a = 0.\]
\end{property}

\subsubsection{Gruppi abeliani}

La struttura algebrica $(\mathbb{Z}, +)$ risulta quindi un gruppo abeliano.
Una generica struttura algebrica $(\mathbb{G}, \cdot)$, formata da un insieme
$\mathbb{G} \neq \emptyset$ e un'operazione binaria "$\cdot$", è un gruppo se e solo se
soddisfa le seguenti proprietà.

\begin{property}
    \textbf{Chiusura di $\mathbb{G}$ su $\cdot$}
    \[\forall a, b \in \mathbb{G}, \; (a \cdot b) \in \mathbb{G}.\]
\end{property}
\begin{property}
    \textbf{Associativa}
    \[\forall a, b, c \in \mathbb{G}, \; (a \cdot b) \cdot c = a \cdot (b \cdot c).\]
\end{property}
\begin{property}
    \textbf{Esistenza dell'elemento neutro}
    \[\forall a \in \mathbb{G}, \; \exists e \in \mathbb{G} \mid a \cdot e = e \cdot a = a.\]
\end{property}
\begin{property}
    \textbf{Esistenza dell'elemento inverso}
    \[\forall a \in \mathbb{G}, \; \exists b \in \mathbb{G} \mid a \cdot b = b \cdot a = e.\]
\end{property}

\begin{observation}
    Non è necessaria la commutatività dell'operazione binaria per definire un gruppo
\end{observation}

Un gruppo si dice inoltre commutativo o abeliano se e solo se soddisfa la seguente proprietà.
\begin{property}
    \textbf{Commutativa}
    \[\forall a, b \in \mathbb{G}, \; a \cdot b = b \cdot a.\]
\end{property}

Affinché una generica struttura algebrica sia un gruppo abeliano,
è dunque necessario che soddisfi contemporaneamente tutte le proprietà sopra elenecate.

\subsubsection{Operazioni in $\mathbb{Z}$}
In $\mathbb{Z}$ è possibile definire l'operazione di prodotto "$\cdot$" nel seguente modo

\[
    \forall a, b \in \mathbb{Z}, \quad a \cdot b := \sum_{k=1}^{b} a.
\]

Si può dimostrare che l'operazione di prodotto gode delle seguenti proprietà.

\begin{property}
    \textbf{Commutativa}
    \[\forall a, b \in \mathbb{Z}, \;a \cdot b = b \cdot a.\]
\end{property}
\begin{property}
    \textbf{Associativa}
    \[\forall a, b, c \in \mathbb{Z}, \;(a \cdot b) \cdot c = a \cdot (b \cdot c).\]
\end{property}
\begin{property}
    \textbf{Esistenza dell'elemento neutro}
    \[\forall a \in \mathbb{Z}, \; a \cdot 1 = 1 \cdot a = a.\]
\end{property}
\begin{property}
    \textbf{Chiusura di $\mathbb{Z}$ rispetto al prodotto}
    \[\forall a, b \in \mathbb{Z}, \; (a \cdot b) \in \mathbb{Z}.\]
\end{property}

\begin{observation}
    All'interno di $\mathbb{Z}$ non può essere definito l'elemento inverso rispetto al prodotto.
\end{observation}
\begin{observation}
    L'operazione di prodotto può essere applicata anche all'insieme $\mathbb{N}_0$, nei confronti della quale anche quest'ultimo risulta chiuso.
\end{observation}



\subsection{Numeri razionali}
Dall'estensione di $\mathbb{Z}$ per includere gli elementi inversi rispetto al prodotto,
si ottiene l'insieme dei numeri razionali, denotato con $\mathbb{Q}$.
Questo insieme gode delle stesse proprietà di $\mathbb{Z}$ rispetto al prodotto,
con l'aggiunta della proprietà dell'esistenza dell'elemento inverso.

L'insieme $\mathbb{Q}$ è definito come l'insieme delle classi di equivalenza individuate
sull'insieme $(\mathbb{Z} \times \mathbb{N})$ dalla seguente relazione.
Definita quindi la relazione $\sim$ nel seguente modo:

\[
    \forall (a, b), (c, d) \in \mathbb{Z} \times \mathbb{N}, \quad (a, b) \sim (c, d) \iff ad = bc,
\]
l'insieme $\mathbb{Q}$ viene definito come

\begin{definition}
    \textbf{Insieme dei numeri razionali}
\[
    \mathbb{Q} := \{p \in \mathbb{Z}, q \in \mathbb{N} \mid \frac{p}{q}\}_{/\sim}
\]
\end{definition}

Ogni elemento di $\mathbb{Q}$ è dunque una classe di equivalenza denotabile come
\[
    \left[r\right]_\sim \; \; con \; \; r \in (\mathbb{Z} \times \mathbb{N})
\]
per comodità di notazione
\[
    r = \left[ \left( p , q \right) \right]_\sim \in \mathbb{Q}
\]
viene indicato con 
\[
    \frac{p_0}{q_0} \mid \frac{p_0}{q_0} \in \left[ \left( p , q \right) \right]_\sim  \land MCD(p_0, q_0) = 1
\]

\subsubsection{Operazioni in $\mathbb{Q}$}
In $\mathbb{Q}$ sono definite le operazioni di somma e prodotto nel seguente modo.

\begin{definition}
    \textbf{Somma in $\mathbb{Q}$}
    \[
        \forall r_1 = \left[\frac{a}{b}\right]_\sim, r_2 = \left[\frac{c}{d}\right]_\sim \in \mathbb{Q}, \;
        r_1 + r_2 := \left[\frac{ad+bc}{bd}\right]
    \]
\end{definition}
\begin{definition}
    \textbf{Prodotto in $\mathbb{Q}$}
    \[
        \forall r_1 = \left[\frac{a}{b}\right]_\sim, r_2 = \left[\frac{c}{d}\right]_\sim \in \mathbb{Q}, \;
        r_1 \cdot r_2 := \left[\frac{ac}{bd}\right]
    \]
\end{definition}

L'operazione di somma in $\mathbb{Q}$ gode delle seguenti proprietà.

\begin{property}
    \textbf{Commutativa}
    \[\forall a, b \in \mathbb{Q}, \, a + b = b + a.\]
\end{property}
\begin{property}
    \textbf{Associativa}
    \[\forall a, b, c \in \mathbb{Q}, \, (a + b) + c = a + (b + c).\]
\end{property}
\begin{property}
    \textbf{Esistenza dell'elemento neutro}
    \[\forall a \in \mathbb{Q}, \, a + 0 = 0 + a = a.\]
\end{property}
\begin{property}
    \textbf{Esistenza dell'elemento inverso}
    \[\forall a \in \mathbb{Q}, \, \exists b = -a \in \mathbb{Q} \mid a + b = 0.\]
\end{property}
\begin{property}
    \textbf{Chiusura di $\mathbb{Q}$ rispetto alla somma}
    \[\forall a, b \in \mathbb{Q}, \; (a + b) \in \mathbb{Q}.\]
\end{property}
La struttura algebrica $(\mathbb{Q}, +)$ è dunque un gruppo abeliano.
\vspace{0.4cm}

L'operazione di prodotto in $\mathbb{Q}' := \mathbb{Q} \setminus \left\{0\right\}$ gode delle seguenti proprietà.

\begin{property}
    \textbf{Commutativa}
    \[\forall a, b \in \mathbb{Q}', \, a \cdot b = b \cdot a.\]
\end{property}
\begin{property}
    \textbf{Associativa}
    \[\forall a, b, c \in \mathbb{Q}', \, (a \cdot b) \cdot c = a \cdot (b \cdot c).\]
\end{property}
\begin{property}
    \textbf{Esistenza dell'elemento neutro}
    \[\forall a \in \mathbb{Q}', \, a \cdot 1 = 1 \cdot a = a.\]
\end{property}
\begin{property}
    \label{p:inv-q}
    \textbf{Esistenza dell'elemento inverso}
    \[\forall a \in \mathbb{Q}', \, \exists b = a^{-1} \in \mathbb{Q}' \mid a \cdot a^{-1} = 1.\]
\end{property}
\begin{property}
    \textbf{Chiusura di $\mathbb{Q}'$ rispetto al prodotto}
    \[\forall a, b \in \mathbb{Q}', \; (a \cdot b) \in \mathbb{Q}.\]
\end{property}
La struttura algebrica $(\mathbb{Q} \setminus \left\{0\right\}, \cdot)$ è dunque un gruppo abeliano.

\begin{observation}
    Viene escluso l'elemento 0 poiché altrimenti non sarebbe garantita \hyperref[p:inv-q]{\textbf{p~\ref*{p:inv-q}}}
\end{observation}


\vspace{0.4cm}

In $\mathbb{Q}$ vale inoltre la proprietà distributiva del prodotto sulla somma
\begin{property}
    \textbf{Distribtiva di $\cdot$ su +}
    \[\forall a, b, c \in \mathbb{Q}, \; a \cdot (b + c) = a \cdot b + a \cdot c.\]
\end{property}

\subsubsection{Relazioni d'ordine totale o lineare}

Una relazione binaria $\sim$ su un insieme $\mathbb{X}$ è detta d'ordine se e solo se valgono le seguenti proprietà.

\begin{property}
    \textbf{Riflessiva}
    \[
        \forall x \in \mathbb{X}, \; x \sim x.
    \]
\end{property}
\begin{property}
    \textbf{Antisimmetrica}
    \[
        x \sim y \land y \sim x \rightarrow x = y.
    \]
\end{property}
\begin{property}
    \textbf{Transitiva}
    \[
        x \sim y \land y \sim z \rightarrow x \sim z.
    \]
\end{property}

\begin{observation}
    Non è richiesto che presi due elementi casuali dell'insieme
    questi possano sempre essere messi in relazione, ma nel momento in cui possano esserlo
    devono valere necessariamente tutte le proprietà di cui sopra.
\end{observation}

\begin{definition}
    \textbf{Relazione d'ordine totale}
    \center

    Una relazione d'ordine generica $\sim$ (d'ordine parziale) definita su $\mathbb{X}$ è
    detta d'ordine totale se e solo se ogni coppia dei elementi dell'insieme può essere
    messa in relazione; in simboli, se viene soddisfatta la seguente proprietà.

    
    \begin{property}
        \textbf{Linearità dell'ordinamento}

        \[ \forall x,y \in \mathbb{X}, \; x \sim y \lor y \sim x \]
    \end{property}
\end{definition}

Prendiamo come relazione su $\mathbb{Q}$, $\leq$.
Si dimostra quindi facilmente che la relazione scelta goda delle proprietà necessarie
per risultare una relazione d'ordine totale.

\begin{proposition}
    \label{rel-ord-tot}
    La relazione $\leq$, definita sugli elementi di $\mathbb{Q}$, è una relazione d'ordine totale.
\end{proposition}

\begin{corollary}
    \textit{di \hyperref[rel-ord-tot]{Pro~\ref*{rel-ord-tot}}}

    \center

    \textbf{$(\mathbb{Q}, +, \cdot, \leq)$ è un campo numerico totalmente ordinato}
\end{corollary}

\end{document}